\subsection{Семинар «Корутины» (6 часов; подведение итогов по Kotlin)}

Для работы с консольным интерфейсом необходимо подключить библиотеку Lanterna. Добавьте следующую зависимость в ваш \texttt{build.gradle.kts} (или аналогичный файл сборки):

\begin{verbatim}
implementation("com.googlecode.lanterna:lanterna:3.1.2")
\end{verbatim}

Библиотека Lanterna позволяет управлять текстовым терминалом: позиционировать курсор, изменять цвета, читать ввод с клавиатуры и очищать экран. Важно помнить, что взаимодействие с терминалом должно происходить из одного потока выполнения, либо вы должны явно обеспечивать потокобезопасность (например, через синхронизацию или использование каналов корутин).

Ниже приведён минимальный шаблон инициализации терминала, совместимый с Windows, macOS и Linux:

\begin{verbatim}
val terminalFactory = DefaultTerminalFactory()
terminalFactory.setForceTextTerminal(false)
terminalFactory.setPreferTerminalEmulator(true)
val terminal = terminalFactory.createTerminal()
terminal.enterPrivateMode()
\end{verbatim}

После выполнения программы обязательно вызовите:

\begin{verbatim}
terminal.exitPrivateMode()
\end{verbatim}

Это восстановит исходное состояние терминала.

Примеры базовых операций:

\begin{itemize}
  \item Установка позиции курсора и вывод строки:
  \begin{verbatim}
terminal.cursorPosition = TerminalPosition(10, 2)
terminal.putString("Hello, World!")
  \end{verbatim}
  
  \item Очистка экрана:
  \begin{verbatim}
terminal.clearScreen()
  \end{verbatim}
  
  \item Чтение ввода:
  \begin{verbatim}
val key = terminal.readInput()
  \end{verbatim}
  
  Поле \verb|key.keyType| содержит тип события (например, \verb|KeyType.Escape|, \verb|KeyType.Enter| и т.д.), а \verb|key.character| — печатаемый символ (если применимо).
\end{itemize}

Все задания ниже должны быть реализованы с использованием корутин Kotlin (\texttt{kotlinx.coroutines}), Flow/Channel/StateFlow/SharedFlow по необходимости, и библиотеки Lanterna для консольного интерфейса. Интерфейс должен быть удобным, информативным и обновляться в реальном времени без миганий или артефактов.

Если какие-то символы сложно вывести, то вы можете заменить их на другие или рисовать с помощью ASCII ART.

\begin{enumerate}
\item[1]
\textbf{Многопользовательский Pomodoro-трекер.} Реализуйте приложение, позволяющее одновременно управлять несколькими Pomodoro-сессиями (например, для разных задач). Каждая сессия — отдельная корутина, генерирующая \texttt{Flow} или отправляющая события в \texttt{Channel} с текущим состоянием: \textit{работа}, \textit{отдых}, \textit{пауза}, \textit{остановлена}, а также оставшимся временем в секундах. Интерфейс должен отображать список всех сессий с возможностью запуска, паузы и остановки каждой из них. Поддержите стандартные интервалы (25/5 мин) и возможность их настройки.

\item[2]
\textbf{Мульти-секундомер для спорта.} Пользователь может создавать именованные секундомеры (например, «Иван», «Мария»). Каждый секундомер поддерживает: запуск, паузу, сброс и фиксацию кругов (laps). Все активные секундомеры отображаются в виде таблицы: имя, текущее время (с точностью до сотых секунды), количество кругов. Время должно обновляться в реальном времени. Каждый секундомер реализуется как независимая корутина, передающая обновления через \texttt{StateFlow} или \texttt{Channel}.

\item[3]
\textbf{Трекер параллельных задач с живым временем.} Пользователь добавляет задачи с названием и длительностью (в секундах). При запуске каждая задача становится корутиной, имитирующей выполнение: прогресс равномерно увеличивается от 0\% до 100\% за заданное время. Пользователь может приостанавливать, возобновлять или отменять любую задачу. Интерфейс отображает таблицу: название, статус (\textit{ожидает}, \textit{выполняется}, \textit{приостановлена}, \textit{завершена}), прогресс (в процентах и визуально, например, через прогресс-бар из символов), прошедшее время. Обновления передаются через \texttt{SharedFlow} или \texttt{StateFlow}.

\item[4]
\textbf{Чат с несколькими агентами.} Пользователь вводит команды вида \texttt{/погода}, \texttt{/курс}, \texttt{/шутка}. Каждая команда обрабатывается отдельной корутиной-агентом, имитирующей сетевой запрос (задержка 1–2 сек). Агенты возвращают текстовые ответы (например, «Сегодня +15°C», «USD = 92.5», «Почему программисты не ходят в лес? Там деревья!»). Все ответы поступают в общий поток и отображаются в хронологическом порядке в нижней части экрана (история сообщений). Верхняя часть содержит поле ввода. Взаимодействие между агентами и UI должно происходить через \texttt{Channel}.

\item[5]
\textbf{Анимированный Pomodoro-трекер с визуальным пульсом.} 
При активной сессии рядом с таймером пульсирует символ: \texttt{$\bullet$} → \texttt{$\circ$} → \texttt{$\bullet$} с периодом 1 сек (500 мс на состояние). Анимация реализуется в отдельной корутине и синхронизируется с основной логикой таймера (например, при паузе анимация останавливается). Интерфейс отображает: текущий режим (работа/отдых), оставшееся время (мин:сек), анимированный индикатор. Все компоненты обмениваются данными через \texttt{Channel}.

\item[6]
\textbf{Таймер для интервальных тренировок (HIIT).} Пользователь задаёт последовательность фаз: название и длительность (например, «Бег — 40 сек», «Отдых — 20 сек»). После запуска приложение циклически проходит фазы. Каждая фаза управляется корутиной, которая отсчитывает время и отправляет обновления. Интерфейс отображает: текущую фазу, оставшееся время, номер цикла (если повторяется), общий прогресс (например, 2/5 циклов завершено). Поддерживается пауза и досрочное завершение. Обновления передаются через \texttt{Flow}.

\item[7]
\textbf{Симулятор нескольких печатающих машинок.} Пользователь создаёт машинки, задавая имя и текст. Каждая машинка посимвольно «печатает» свой текст в выделенной строке с задержкой (например, 100 мс на символ). Пользователь может остановить или возобновить печать любой машинки. Каждая машинка — отдельная корутина, отправляющая своё текущее состояние (напечатанная часть текста, позиция курсора) в UI через \texttt{Channel}. Интерфейс отображает все машинки одновременно.

\item[8]
\textbf{Менеджер фоновых загрузок (симуляция).} Пользователь добавляет загрузку, указывая имя, объём (например, 1000 условных единиц) и скорость (единиц/сек). Каждая загрузка — корутина, которая постепенно увеличивает прогресс. Интерфейс отображает таблицу: имя, прогресс (в процентах и визуально), текущая скорость, оставшееся время (расчётное). Любую загрузку можно отменить. Обновления передаются через \texttt{StateFlow} или \texttt{SharedFlow}.

\item[9]
\textbf{Живой лог-монитор с цветовой маркировкой.} Создайте три источника логов: «Сервер», «База данных», «Фронтенд». Каждый — отдельная корутина, генерирующая сообщения с разной частотой (например, сервер — каждые 500 мс, БД — каждые 1.2 сек, фронтенд — каждые 800 мс). Сообщения поступают в общий \texttt{BroadcastChannel} (или \texttt{SharedFlow}) и отображаются в UI как список последних 15 строк. Сообщения от разных источников визуально различаются: цветом текста или префиксом вида \texttt{[SRV]}, \texttt{[DB]}, \texttt{[FE]}. Пользователь может приостановить любой источник.

\item[10]
\textbf{Трекер времени для нескольких настольных игр за вечер.} Пользователь может запускать таймеры для разных игр («Шахматы», «Каркассон» и т.д.). Каждая игра может быть активна независимо от других. Для каждой игры отображается: название, статус (активна/приостановлена), общее потраченное время. Каждая игра управляется отдельной корутиной, передающей обновления через \texttt{Flow}. Поддержите одновременный запуск нескольких активных игр.

\item[11]
\textbf{Симулятор частиц огня.} Пользователь нажимает клавишу (например, \texttt{SPACE}), чтобы запустить новую частицу (ASCII-символ: \texttt{*}, \texttt{o}, \texttt{+}). Каждая частица движется вверх по экрану с индивидуальной скоростью (например, задержка 100–300 мс между шагами). При выходе за верхнюю границу частица удаляется. Пользователь может остановить все частицы. Каждая частица — отдельная корутина, отправляющая свою текущую позицию в UI через \texttt{Channel}. Интерфейс перерисовывает все частицы на каждом кадре.

\item[12]
\textbf{Симулятор нескольких аквариумов.} Каждый аквариум — отдельная строка, содержащая одну «рыбку» (например, \texttt{><>}). Рыбка плавает вправо-влево в пределах своей строки с постоянной скоростью. Пользователь может: создать аквариум (с именем), удалить существующий, приостановить движение конкретной рыбки. Каждая рыбка — отдельная корутина с циклической анимацией. Интерфейс отображает все аквариумы одновременно. Обновления позиций передаются через \texttt{Flow}.

\item[13]
\textbf{Таймер для языковой практики.} Пользователь запускает цикл из трёх фаз: «Говорение — 2 мин», «Письмо — 1 мин», «Отдых — 30 сек». После завершения фазы автоматически начинается следующая; после последней — цикл повторяется. Интерфейс отображает: текущую фазу, оставшееся время, номер цикла, общий прогресс. Поддерживается пауза и пропуск текущей фазы (переход к следующей). Каждая фаза управляется корутиной, смена фаз — через \texttt{Channel}.

\item[14]
\textbf{Мульти-таймер для готовки по рецепту.} Рецепт состоит из шагов, некоторые из которых помечены как параллельные. Например: «Варить пасту — 10 мин» (параллельно), «Готовить соус — 7 мин» (параллельно), «Смешать — 1 мин» (последовательно). При запуске все параллельные шаги начинают отсчёт одновременно. Интерфейс отображает каждый шаг с оставшимся временем, статусом (\textit{ожидает}, \textit{выполняется}, \textit{завершён}) и индикатором параллельности. Каждый шаг — отдельная корутина, обновления передаются через \texttt{Flow}.

\item[15]
\textbf{Трекер времени в приложениях.} Пользователь может «включить» приложение (например, «Браузер», «IDE»). При включении запускается таймер использования; при выключении — останавливается, но накопленное время сохраняется. Интерфейс отображает список приложений с общим временем за текущую сессию и статусом (активно/неактивно). Каждое приложение управляется отдельной корутиной, которая отслеживает активное время и отправляет обновления через \texttt{StateFlow}.

\item[16]
\textbf{Симулятор нескольких будильников.} Пользователь создаёт будильник, указывая либо абсолютное время срабатывания (ЧЧ:ММ), либо интервал (например, каждые 5 минут). При срабатывании в UI появляется уведомление (например, мигающая строка). Пользователь может: «отложить» на 1 минуту или отключить. Каждый будильник — отдельная корутина, ожидающая события через \texttt{delay} или \texttt{withTimeout}. Уведомления передаются в UI через \texttt{Channel}.

\item[17]
\textbf{Трекер велопоездок с точками останова.} Поездка состоит из последовательности отрезков: «Дом → Парк — 15 мин», «Парк → Река — 20 мин» и т.д. После завершения отрезка поездка автоматически переходит в режим паузы. Пользователь может возобновить движение (начать следующий отрезок) или завершить поездку досрочно. Интерфейс отображает: текущий отрезок, общее пройденное время, статус (\textit{в пути}, \textit{пауза}, \textit{завершена}). Каждый отрезок управляется корутиной.

\item[18]
\textbf{Живой генератор цитат.} Цитаты из заранее заданного списка появляются в UI одна за другой с интервалом (например, каждые 30 сек). Пользователь может: «сохранить» текущую цитату (она перемещается в отдельный блок истории), пропустить текущую или остановить генерацию. Генерация реализована в отдельной корутине, взаимодействие с UI — через \texttt{StateFlow}, содержащий текущую цитату и историю.

\item[19]
\textbf{Трекер медитаций с дыхательной анимацией.} 
Во время сессии в центре экрана отображается символ (например, \texttt{$\circ$}), который плавно «расширяется» (вдох — 4 сек) и «сужается» (выдох — 6 сек). Анимация реализуется в отдельной корутине (через постепенное изменение символа или его окружения), а таймер сессии — в другой. Обе корутины синхронизированы через \texttt{Channel} (например, для паузы или остановки). Поддерживается пауза и остановка сессии.

\item[20]
\textbf{Pomodoro-трекер с историей и статистикой.} Пользователь запускает Pomodoro-сессии (работа/отдых). После завершения сессия сохраняется в историю (последние 10 записей с временем завершения и типом). Интерфейс разделён на две области: верхняя — активные сессии (с управлением), нижняя — история (только просмотр). Активные сессии реализованы как корутины, завершённые — хранятся в \texttt{SharedFlow} или \texttt{MutableStateFlow}.

\item[21]
\textbf{Симулятор мигающих огней на карте.} Пользователь добавляет «огонёк» в заданную строку экрана (например, вводит номер строки). Каждый огонёк мигает с индивидуальной частотой: \texttt{*} → пробел → \texttt{*} с заданным периодом (например, 800 мс). Каждый — отдельная корутина. Пользователь может удалить любой огонёк. Весь экран перерисовывается на основе актуальных состояний, передаваемых через \texttt{Channel}.

\item[22]
\textbf{Таймер для записи подкаста с главами.} Пользователь запускает запись — идёт общий таймер. В любой момент можно нажать клавишу (например, \texttt{C}), чтобы добавить «главу» — фиксируется временная метка (мин:сек). Интерфейс отображает: общий таймер и список глав с их временными метками. Запись (таймер) и обработка глав — разные корутины, взаимодействие через \texttt{Channel}.

\item[23]
\textbf{Трекер времени на улице.} Пользователь запускает «прогулку» — начинается отсчёт времени. При входе в помещение (магазин, дом) можно поставить на паузу; при выходе — возобновить. Можно вести несколько таких сессий в течение дня (например, «Утро», «Вечер»). Каждая сессия — отдельная корутина с состоянием (\textit{активна}/\textit{пауза}) и накопленным временем. UI отображает все сессии с возможностью управления.

\item[24]
\textbf{Симулятор нескольких сердец с разным ритмом.} Пользователь создаёт «сердце», 
задавая частоту (ударов в минуту, например, 60, 72, 80). 
Каждое сердце отображается как символ \texttt{\Heart}, который мигает в заданном ритме
 (например, при 60 уд/мин — вспышка каждую секунду). Каждое — отдельная корутина с \texttt{delay}, 
 рассчитанным по формуле $ \text{delay} = 60\,000 / \text{частота} $ мс. 
 Пользователь может остановить любое сердце. Анимация обновляется через \texttt{Flow}.

\item[25]
\textbf{Таймер хода для настольных игр.} Пользователь задаёт список игроков и время на ход (например, 60 сек). После старта таймер идёт для текущего игрока. По истечении времени — звуковое или визуальное уведомление и автоматический переход к следующему игроку. Пользователь может продлить ход (добавить 30 сек) или пропустить игрока. Каждый игрок управляется логикой в корутине, обновления состояния (текущий игрок, оставшееся время) передаются через \texttt{StateFlow}.

\item[26]
\textbf{Мульти-таймер для йоги с визуальной подсказкой.} Пользователь выбирает последовательность поз (например, из списка: «Гора», «Дерево», «Собака мордой вниз»). Каждая поза отображается как упрощённая ASCII-схема (например, \texttt{/|\textbackslash{}} для «горы»). Для каждой позы идёт таймер (например, 30 сек). После окончания можно: автоматически перейти к следующей или ждать подтверждения пользователя (нажатие клавиши). Каждая поза — корутина, UI обновляется по \texttt{Flow}.

\item[27]
\textbf{Консольный космос: спутники на орбите.} Пользователь добавляет спутник, задавая радиус орбиты (в символах) и угловую скорость (градусов/сек). Каждый спутник движется по круговой траектории в пределах ASCII-экрана с использованием упрощённых тригонометрических расчётов (например, \texttt{x = centerX + radius * cos(angle)}, \texttt{y = centerY + radius * sin(angle)}). Каждый спутник — отдельная корутина с \texttt{delay}, обновляющая свою позицию. Пользователь может удалить спутник. Положения всех спутников передаются в UI через \texttt{Channel} для перерисовки.

\item[28]
\textbf{Менеджер Pomodoro-сессий для команды.} Пользователь вводит имена участников (например, «Анна», «Борис»). Для каждого можно запустить независимую Pomodoro-сессию. Интерфейс отображает таблицу: имя, статус (\textit{работа}/\textit{отдых}/\textit{пауза}), оставшееся время. Управление осуществляется клавишами: стрелки — выбор участника, \texttt{Enter} — запуск/пауза, \texttt{Del} — остановка. Каждая сессия — корутина, обновления — через \texttt{SharedFlow}.
\end{enumerate}